% GENERATED FILE. Edit canonical lessons, then run npm run generate:books.
% canonical-source-hash: fnv1a64:bdf9a3db509ea4dd
% canonical-lessons: HI-C136-kurta, HI-C136-sari, HI-C136-dupatta, HI-C136-topi, HI-C136-kot

\chapter{Kurta, Sari, Scarf, Cap, Coat}
\label{ch:hi-kurta-sari-scarf-cap-coat}
\hlchaptermodality{136}

\begin{chapteropening}
\textbf{By the end of this chapter:} \emph{I can say a kurta, a sari, a scarf, a cap and a coat.}

Complete the last lesson of chapter 136: I can say a kurta, a sari, a scarf, a cap and a coat.
\end{chapteropening}

\section[\texorpdfstring{\hi{कुर्ता} (kurtā)}{कुर्ता (kurtā)}]{\hi{कुर्ता} (kurtā) — a kurta}

\label{lesson:HI-C136-kurta}



\begin{quote}
Before the new one: say the Hindi for a shirt, then the Hindi for trousers.

Say \textbf{\hi{हरा}} and \textbf{\hi{पीला}}, green and yellow, two colour words Hindi inherited.
\end{quote}

\subsection*{You'll want to know: \hi{कुर्ता}}
\textbf{\hi{कुर्ता}} — \emph{kurtā} — \textquotedblleft{}a kurta\textquotedblright{}.

\begin{cousinweb}
\textbf{\hi{कुर्ता}} (\emph{kurtā}) is a loan, usually traced to Persian \emph{kurta}, \textquotedblleft{}a long shirt\textquotedblright{}.
\end{cousinweb}

\begin{grammarlens}[title={What you've built}]
One more word for this chapter.
\end{grammarlens}

\subsection*{Your turn}
\begin{itemize}
\raggedright
  \item \emph{Say it:} \emph{kurtā}
  \item \emph{Say it:} \emph{kurtā}, once more
  \item \emph{Say it:} say \emph{kurtā}
  \item \emph{Recall:} say the Hindi for a shirt, then the Hindi for trousers, then say \emph{kurtā} again
\end{itemize}

\begin{tcolorbox}[breakable,colback=teal!4,colframe=teal!35!black,title={Before you move on}]
What does \hi{कुर्ता} mean? (\textquotedblleft{}A kurta\textquotedblright{}.) Say it once more.
\end{tcolorbox}
\section[\texorpdfstring{\hi{साड़ी} (sāṛī)}{साड़ी (sāṛī)}]{\hi{साड़ी} (sāṛī) — a sari}

\label{lesson:HI-C136-sari}



\begin{quote}
Before the new one: say the Hindi for trousers, then the Hindi for a kurta.

Say the colours and their stories: \textbf{\hi{लाल}}, a red gem's name; \textbf{\hi{नीला}}, indigo's dark blue; \textbf{\hi{सफ़ेद}}, a Persian loan.
\end{quote}

\subsection*{You'll want to know: \hi{साड़ी}}
\textbf{\hi{साड़ी}} — \emph{sāṛī} — \textquotedblleft{}a sari\textquotedblright{}.

\begin{cousinweb}
\textbf{\hi{साड़ी}} (\emph{sāṛī}) is a tadbhava, a word inherited from Sanskrit and reshaped by sound change. It comes from \emph{śāṭī}, \textquotedblleft{}strip of cloth, garment\textquotedblright{} (CDIAL 12381). CDIAL notes that the word also travelled west, to Pashto and to Persian \emph{shāl}, the source of English \emph{shawl}.
\end{cousinweb}

\begin{grammarlens}[title={What you've built}]
One more word for this chapter.
\end{grammarlens}

\subsection*{Your turn}
\begin{itemize}
\raggedright
  \item \emph{Say it:} \emph{sāṛī}
  \item \emph{Say it:} \emph{sāṛī}, once more
  \item \emph{Say it:} say \emph{sāṛī}
  \item \emph{Recall:} say the Hindi for trousers, then the Hindi for a kurta, then say \emph{sāṛī} again
\end{itemize}

\begin{tcolorbox}[breakable,colback=teal!4,colframe=teal!35!black,title={Before you move on}]
What does \hi{साड़ी} mean? (\textquotedblleft{}A sari\textquotedblright{}.) Say it once more.
\end{tcolorbox}
\section[\texorpdfstring{\hi{दुपट्टा} (dupaṭṭā)}{दुपट्टा (dupaṭṭā)}]{\hi{दुपट्टा} (dupaṭṭā) — a scarf, a dupatta}

\label{lesson:HI-C136-dupatta}



\begin{quote}
Before the new one: say the Hindi for a kurta, then the Hindi for a sari.
\end{quote}

\subsection*{You'll want to know: \hi{दुपट्टा}}
\textbf{\hi{दुपट्टा}} — \emph{dupaṭṭā} — \textquotedblleft{}a scarf, a dupatta\textquotedblright{}.

\begin{cousinweb}
\textbf{\hi{दुपट्टा}} (\emph{dupaṭṭā}) is built from \emph{du-}, \textquotedblleft{}two\textquotedblright{}, and \emph{paṭṭā}, \textquotedblleft{}a strip of cloth\textquotedblright{}: a cloth of two breadths. CDIAL lists Hindi \emph{paṭṭī}, \textquotedblleft{}strip of cloth\textquotedblright{}, under Sanskrit \emph{paṭṭa}, \textquotedblleft{}cloth\textquotedblright{}, probably of non-Aryan source (CDIAL 7700), but not \emph{dupaṭṭā} itself.
\end{cousinweb}

\begin{grammarlens}[title={What you've built}]
One more word for this chapter.
\end{grammarlens}

\subsection*{Your turn}
\begin{itemize}
\raggedright
  \item \emph{Say it:} \emph{dupaṭṭā}
  \item \emph{Say it:} \emph{dupaṭṭā}, once more
  \item \emph{Say it:} say \emph{dupaṭṭā}
  \item \emph{Recall:} say the Hindi for a kurta, then the Hindi for a sari, then say \emph{dupaṭṭā} again
\end{itemize}

\begin{tcolorbox}[breakable,colback=teal!4,colframe=teal!35!black,title={Before you move on}]
What does \hi{दुपट्टा} mean? (\textquotedblleft{}A scarf, a dupatta\textquotedblright{}.) Say it once more.
\end{tcolorbox}
\section[\texorpdfstring{\hi{टोपी} (ṭopī)}{टोपी (ṭopī)}]{\hi{टोपी} (ṭopī) — a cap, a hat}

\label{lesson:HI-C136-topi}



\begin{quote}
Before the new one: say the Hindi for a sari, then the Hindi for a scarf.
\end{quote}

\subsection*{You'll want to know: \hi{टोपी}}
\textbf{\hi{टोपी}} — \emph{ṭopī} — \textquotedblleft{}a cap, a hat\textquotedblright{}.

\begin{cousinweb}
\textbf{\hi{टोपी}} (\emph{ṭopī}) is a tadbhava, a word inherited and reshaped by sound change. It comes from \emph{*ṭoppa}, \textquotedblleft{}hat, covering\textquotedblright{}, a form CDIAL reconstructs from its descendants (CDIAL 5481). CDIAL calls any link with a Munda source very doubtful.
\end{cousinweb}

\begin{grammarlens}[title={What you've built}]
One more word for this chapter.
\end{grammarlens}

\subsection*{Your turn}
\begin{itemize}
\raggedright
  \item \emph{Say it:} \emph{ṭopī}
  \item \emph{Say it:} \emph{ṭopī}, once more
  \item \emph{Say it:} say \emph{ṭopī}
  \item \emph{Recall:} say the Hindi for a sari, then the Hindi for a scarf, then say \emph{ṭopī} again
\end{itemize}

\begin{tcolorbox}[breakable,colback=teal!4,colframe=teal!35!black,title={Before you move on}]
What does \hi{टोपी} mean? (\textquotedblleft{}A cap, a hat\textquotedblright{}.) Say it once more.
\end{tcolorbox}
\section[\texorpdfstring{\hi{कोट} (koṭ)}{कोट (koṭ)}]{\hi{कोट} (koṭ) — a coat}

\label{lesson:HI-C136-kot}



\begin{quote}
Before the new one: say the Hindi for a scarf, then the Hindi for a cap.
\end{quote}

\subsection*{You'll want to know: \hi{कोट}}
\textbf{\hi{कोट}} — \emph{koṭ} — \textquotedblleft{}a coat\textquotedblright{}.

\begin{cousinweb}
\textbf{\hi{कोट}} (\emph{koṭ}) is a loan from English \emph{coat}. It is a different word from Hindi \emph{koṭ}, \textquotedblleft{}fort\textquotedblright{}, from Sanskrit \emph{koṭṭa} (CDIAL 3500).
\end{cousinweb}

\begin{grammarlens}[title={What you've built}]
One more word for this chapter.
\end{grammarlens}

\subsection*{Your turn}
\begin{itemize}
\raggedright
  \item \emph{Say it:} \emph{koṭ}
  \item \emph{Say it:} \emph{koṭ}, once more
  \item \emph{Say it:} say \emph{koṭ}
  \item \emph{Recall:} say the Hindi for a scarf, then the Hindi for a cap, then say \emph{koṭ} again
\end{itemize}

\begin{tcolorbox}[breakable,colback=teal!4,colframe=teal!35!black,title={Before you move on}]
What does \hi{कोट} mean? (\textquotedblleft{}A coat\textquotedblright{}.) Say it once more.
\end{tcolorbox}
